\begin{prerequis}
Lire un repère, calculer la valeur d'une expression, résoudre graphiquement une
équation simple et distinguer abscisse et ordonnée.
\end{prerequis}
\begin{automatismes}
Pour $f(x)=x^2-3$, calculer $f(-2)$, $f(0)$ et $f(3)$; résoudre
$f(x)=1$; expliquer pourquoi $f(-2)=1$ ne signifie pas $f(1)=-2$.
\end{automatismes}
\begin{activite}[Activité d'entrée -- Une machine, plusieurs représentations]
Une machine associe à $x$ le nombre $x^2-4$. Construire une table pour sept
valeurs de $x$, placer les points correspondants, puis répondre aux mêmes
questions à partir de la formule, du tableau et du graphique. Comparer ce que
chaque représentation rend facile ou difficile.
\end{activite}

